\chapter{Conclusões e trabalhos futuros}
\label{cap:conclusoes}

A comparação das especificações 0126 e 1580 mostrou que, para o Edifício Palatium, as alterações de maior repercussão se concentram na disposição dos equipamentos de entrada e na documentação das decisões de dimensionamento. Entre os 82 requisitos examinados para o Ceará, 38 foram classificados como alterados e 17 como acrescentados. A incidência sobre o edifício foi verificada a partir das características da rede e das soluções representadas no projeto. Esse procedimento selecionou 39 requisitos, correspondentes a 38 registros de impacto. A matriz permite recuperar, para cada decisão, o dispositivo normativo, a evidência do projeto e as condições que limitam o resultado.

O posicionamento do centro de proteção geral junto ao alinhamento da propriedade altera o percurso dos condutores e exige compatibilização com a infraestrutura civil. O cenário estudado reúne cinco frentes físicas: reposicionar o centro de proteção geral, prever cubículo exclusivo para o centro de medição, deslocar a caixa de medição de incêndio, prever a proteção contra surtos junto à entrada e adequar a identificação e o acesso aos conjuntos. A medição de incêndio ilustra a distinção entre uma mudança de localização e uma mudança de circuito. O projeto já apresenta a derivação a montante do disjuntor geral. Com a caixa em sua nova posição, permanecem 28 medições, sendo 27 no centro convencional e uma independente.

A comparação da demanda revelou uma limitação da documentação disponível. A resposta da Enel esclareceu que o método da especificação~0126 pode servir de referência técnica em conjunto com a NBR~5410, desde que a memória explicite a metodologia e os fatores utilizados. O cálculo reproduziu os 291,33~kVA declarados no memorial. A origem do multiplicador residencial de 3,2 e a contribuição simultânea dos carregadores permanecem indeterminadas. Assim, os 96,20~kW de carga instalada de recarga identificados na reconstituição não permitem estabelecer uma nova demanda coletiva. Também dependem de dados complementares a capacidade de interrupção requerida no ponto de instalação e as decisões vinculadas à disponibilidade da rede.

Os resultados descrevem o cenário de aplicação da especificação~1580 ao projeto aprovado, com o alcance próprio de um estudo documental. Como continuidade, propõe-se confrontar a demanda calculada com medições da entrada geral, registrando a ocupação dos apartamentos e os períodos de recarga. Outra etapa consiste em verificar o percurso até o novo centro de proteção geral e o espaço de manutenção do centro de medição em conjunto com os projetos de arquitetura, estrutura e drenagem. Essas verificações permitiriam passar das consequências identificadas à avaliação das alternativas de implantação.
